\section{Discussion, Limitations, and Research Agenda}
\label{sec:discussion}
\subsection{What the Evidence Establishes}
The saved artifacts show consistency across three representations: a requested
network action, the accepted operation-level schedule, and the resulting
electrical measurements. The isolation-only and physical-fault cases separate
two phenomena that a purely graphical demonstration might conflate.
Disconnecting a bus is not itself a short circuit; the native fault case
adds an independently scheduled conductive path and measured fault current.
Similarly, tie closure is supported by measured branch current and downstream
voltage recovery, not just a changed line color.

The workflow is useful as a collaboration boundary. Networking colleagues can
provide candidate routes without becoming responsible for PSCAD component
identifiers, while electrical-model development can proceed without rewriting
their routing code. The operation and measurement records make disagreements
inspectable: a missing graph path, rejected closure, invalid mapping, failed
build, stale output, and an electrically poor result are different failure
classes. This separation supports debugging and subsequent experimental design.

Nevertheless, traceability is not equivalent to electrical admissibility.
Radiality and source connectivity are necessary modeling checks for the chosen
scenario, but they are not a complete operating envelope. An overload or
unacceptable voltage may occur in an acyclic, source-connected network. A
future controller must use calibrated electrical constraints and independent
interlocks before these commands could be considered for hardware.

\subsection{Threats to Validity}
\paragraph{Electrical validity.}
The six-bus network is intentionally small and uses illustrative branch
impedances, source settings, and resistive loads. The load-sizing scalar is
not the load power factor, and demand is voltage dependent. There is no
validated transformer, grounding arrangement, DER controller, switching arc,
protection curve, or grid-forming island. The restored voltage range should
not be generalized to a campus or community microgrid without a new parameter
set and validation.

\paragraph{Control and safety semantics.}
The implemented route checks are local and operation based. Partial command
execution is possible, and command warnings are not equivalent to a
transactional abort. This is a material limitation for generalizing the
loop-rejection result. Atomic validation, explicit safe intermediate states,
independent permissives, and rollback or recovery policies should precede
any live or hardware-facing controller integration.

\paragraph{Timing and communications.}
No independent communication failure is simulated in the electrical run.
The control delay is an input, not an observed transport characteristic.
An always-connected control-plane objective remains untested, including its
power supply, redundancy, reconnection behavior, and failure isolation.
The MCP service is a local automation mechanism, not a safety-rated
field protocol or evidence of secure, reliable network delivery.

\paragraph{Measurement and repeatability.}
The analysis threshold and interval filtering are experimental definitions.
The selected breaker-state traces are controller outputs rather than
independent contacts. The stored 501-point previews cannot support detailed
fault-waveform reconstruction. No timestep-convergence study, independently
recomputed per-scenario raw-output analysis, repeated-run variance, or
instrument calibration is supplied with this draft. The saved deterministic
matrix supports functional inspection, not statistical reliability claims.

\paragraph{Novelty and comparison.}
Software-defined microgrids and cyber-physical testbeds already exist. The
current manuscript does not show that its automation architecture outperforms
those systems, or that the networking algorithm is optimal. A future venue
submission must align its contribution with an appropriately scoped systems
or experimental-methods claim, or add a genuinely new and evaluated control
method.

\subsection{Next Experimental Stage}
The next stage should first freeze measured or benchmark electrical
parameters, validate the steady state and fault level independently, and
repeat the cases at multiple EMT time steps. A larger feeder and unbalanced
loads would test whether the adapter remains useful beyond the current
balanced six-bus example. Branch ampacity, source limits, neutral and grounding
models, and load dynamics should be explicit inputs rather than inferred from
network-capacity labels.

Live integration then needs an agreed controller transport with versioned
messages, identifiers, acknowledgments, stale-command handling, and a
well-defined failure policy. Fault detection must originate from measured
electrical signals if the experiment is to evaluate protection response.
Closing the measurement-to-controller loop also requires an explicit clock
and synchronization strategy; established federation approaches such as
HELICS and FNCS provide relevant design precedents
\cite{palmintier2017,ciraci2014}, but choosing or integrating one is future
work.

A useful comparative protocol would hold the electrical case fixed and vary
only the controller and communication treatment. Candidate conditions are
manual or fixed restoration, the current supplied-route policy, and the
networking team's updated policy. Controlled delay, loss, disconnection, and
recovery distributions should use documented seeds and repeated trials.
Record detection time, issue time, receipt time, breaker operation, and load
recovery separately; do not collapse them into a single ``SDN delay.''

The evaluation should report operation-level and transaction-level rejection,
voltage-limit violations, loading limits, switching count, recovery duration,
and unavailable demand. For example, a future energy-not-served measure could
be defined using a specified counterfactual demand trajectory:
\begin{equation}
 E_{\mathrm{NS}}=\sum_i\int
 \max\{P_i^{\mathrm{required}}(t)-P_i^{\mathrm{served}}(t),0\}\,dt.
\end{equation}
This metric is not computed here. In a resistive-load model, reduced consumed
power during undervoltage is not automatically equivalent to disconnected
customer demand; the counterfactual must be stated before calculating such
an index. Comparative results should include uncertainty appropriate to the
randomized treatments and retain all raw electrical and communication logs.
